\subsection{Transfer to later floods}
\label{sec:temporal}

The held-out blocks separate training from testing in space but not in time, since the labels and the reference extents come from the same events, while a map is meant to anticipate a flood that has not yet happened. Each flood model is therefore also trained on labels from the events up to 2022, which are 2017, 2019 and 2020 in Rangpur and Rajshahi, 2017 and 2022 in Sylhet and the three cyclones from 2020 to 2022 on the coast, and scored on the held-out blocks against the extents of the 2024 events. The training labels keep each region's rule, capped by the number of earlier events, so ground counts as flood-prone where it flooded in at least two of three events in Rangpur and Rajshahi, in both events in Sylhet and in at least one cyclone on the coast, and it counts as flooded in 2024 where any 2024 event reached it. Three references are scored on the same assets over the same \RRTempSeeds{} block assignments. The first is the share of the earlier events in which the ground flooded, which is the map a planner would hold from the radar record alone, the second is the same model trained on labels that include the 2024 events, which is the spatial-only design of Section~\ref{sec:observed}, and the third is the same model trained on the terrain-threshold labels of Section~\ref{sec:labels}, which draw on no flood event.

\begin{table*}[t]
\centering
\small
\begin{tabular}{lcccc}
\toprule
Region & Model to 2022 & All events & Past flooding & Terrain labels \\
\midrule
Rangpur \& Rajshahi & \RRTempAUC{} $\pm$ \RRTempSD{} & \RRTempAllAUC{} & \RRTempPastAUC{} $\pm$ \RRTempPastSD{} & \RRTempProxyAUC{} $\pm$ \RRTempProxySD{} \\
Sylhet & \SylhetTempAUC{} $\pm$ \SylhetTempSD{} & \SylhetTempAllAUC{} & \SylhetTempPastAUC{} $\pm$ \SylhetTempPastSD{} & \SylhetTempProxyAUC{} $\pm$ \SylhetTempProxySD{} \\
South-west coast & \CoastalTempAUC{} $\pm$ \CoastalTempSD{} & \CoastalTempAllAUC{} & \CoastalTempPastAUC{} $\pm$ \CoastalTempPastSD{} & \CoastalTempProxyAUC{} $\pm$ \CoastalTempProxySD{} \\
\bottomrule
\end{tabular}
\caption{AUC against the 2024 flood extents on held-out 10\,km test blocks, as mean $\pm$ standard deviation over \RRTempSeeds{} block assignments. The model to 2022 is trained only on the earlier events, the all-events model on labels that include 2024, past flooding is the share of the earlier events in which the ground flooded, and the terrain-label model is the same model trained on the terrain-threshold labels. The 2024 events reached \RRTempTestRate{}\,\%, \SylhetTempTestRate{}\,\% and \CoastalTempTestRate{}\,\% of assets in the three regions.}
\label{tab:temporal}
\end{table*}

Trained only on the earlier floods, the model ranks the ground the 2024 floods reached at an AUC of \RRTempAUC{} in Rangpur and Rajshahi, \SylhetTempAUC{} in Sylhet and \CoastalTempAUC{} on the coast (Table~\ref{tab:temporal}). In Rangpur and Rajshahi leaving the 2024 events out of training costs nothing measurable ($p\RRTempMinusAllP{}$), while in Sylhet and on the coast it costs \SylhetTempAllGap{} and \CoastalTempAllGap{} in AUC, which is the price of anticipating a flood the model has not seen. The comparison with past flooding is the more instructive one. In both riverine regions the radar record of the earlier floods ranks the 2024 flooding better than the model does, at \RRTempPastAUC{} against \RRTempAUC{} in Rangpur and Rajshahi, where the model is ahead on \RRTempAheadPast{} of \RRTempSeeds{} assignments, and \SylhetTempPastAUC{} against \SylhetTempAUC{} in Sylhet, where it is ahead on \SylhetTempAheadPast{}. On the coast the order reverses, with the model ahead at \CoastalTempAUC{} against \CoastalTempPastAUC{} on \CoastalTempAheadPast{} of \CoastalTempSeeds{} assignments.

The likely reading is that riverine flooding returns to much the same ground from one monsoon to the next, so that the footprint of several observed floods is already a strong guide to the next one, whereas each cyclone surge floods the stretch of coast where it makes landfall. The 2024 floods bear this out, since \RRPastFeatUnseenPos{}\,\% of the flooded assets in Rangpur and Rajshahi and \SylhetPastFeatUnseenPos{}\,\% in Sylhet stood on ground no earlier event had reached, against \CoastalPastFeatUnseenPos{}\,\% on the coast.

The record of earlier flooding can also be given to the model as a feature, provided it comes from floods before the ones the model learns to predict, since a feature drawn from the training events would let the model copy its own labels. The model is therefore trained to predict the flooding of the last event year up to 2022 from its usual features and the share of the events before that year in which the ground flooded, and it is scored on the held-out blocks against the 2024 floods with the share taken over all the events up to 2022. In the riverine regions the combination outperforms both of its parts. It reaches \RRPastFeatAUC{} in Rangpur and Rajshahi and \SylhetPastFeatAUC{} in Sylhet, against \RRPastFeatOnlyAUC{} and \SylhetPastFeatOnlyAUC{} for past flooding alone, ahead on \RRPastFeatOverPastAhead{} and \SylhetPastFeatOverPastAhead{} of \RRTempSeeds{} assignments ($p\RRPastFeatOverPastP{}$ and $p\SylhetPastFeatOverPastP{}$), and against \RRPastFeatWithoutAUC{} and \SylhetPastFeatWithoutAUC{} for the same model trained on the same target without the feature. The terrain and access features therefore add to the flood record rather than repeat it, and they do so on the ground no earlier flood reached, where past flooding gives every asset the same score of zero. That ground holds \RRPastFeatUnseenShare{}\,\% of the test assets in Rangpur and Rajshahi and \SylhetPastFeatUnseenShare{}\,\% in Sylhet, and on it the combined model ranks the 2024 flooding at \RRPastFeatUnseenAUC{} and \SylhetPastFeatUnseenAUC{}, where past flooding alone can do no better than chance. On the coast the design is weaker, because it leaves a single cyclone to train on, and the combination at \CoastalPastFeatAUC{} cannot be told apart from past flooding alone at \CoastalPastFeatOnlyAUC{} ($p\CoastalPastFeatOverPastP{}$), both well below the \CoastalTempAUC{} of the model trained on all three earlier cyclones. The riverine maps published with this study use the feature, trained in this way and applied with the share taken over every mapped event (Section~\ref{sec:gnn}), while every other validation figure comes from the model without it.

The 2024 floods also give the label comparison of Section~\ref{sec:labelresult} a test that favours neither label source. There the model trained on observed extents is scored against the definition it learned and the model trained on terrain labels against one it never saw, whereas neither has seen the 2024 floods. The observed labels still lead, at \RRTempAUC{} against \RRTempProxyAUC{} in Rangpur and Rajshahi, \SylhetTempAUC{} against \SylhetTempProxyAUC{} in Sylhet and \CoastalTempAUC{} against \CoastalTempProxyAUC{} on the coast, ahead on \RRTempProxyAhead{}, \SylhetTempProxyAhead{} and \CoastalTempProxyAhead{} of \RRTempSeeds{} assignments ($p\RRTempProxyGapP{}$, $p\SylhetTempProxyGapP{}$ and $p\CoastalTempProxyGapP{}$). The gap of \RRTempProxyGap{} in Rangpur and Rajshahi matches the \RRLabelGap{} of the spatial comparison and the coastal gap of \CoastalTempProxyGap{} stays close to its \CoastalLabelGap{}, while in Sylhet the gap falls from \SylhetLabelGap{} to \SylhetTempProxyGap{}, so part of the Sylhet advantage in the spatial comparison came from scoring against the definition the observed model was trained on rather than from the labels themselves.
